\chapter{प्रस्तावना}

इस पुस्तक की महत्वाकांक्षा सीधी-सादी है: रसायन विज्ञान की एक ही सुसंगत
पुस्तक, जो विद्यालय के पहले वर्ष से लेकर स्नातक की डिग्री के अंत तक सब कुछ
समेटे, और जो दुनिया में कहीं भी पढ़ने वालों के लिए हो। यह पहले अंग्रेज़ी में
लिखी गई थी; आपके हाथ में उसका हिंदी संस्करण है।

इसकी शैली संक्षिप्त और कसी हुई है: परिभाषाओं, उदाहरणों, प्रतिज्ञप्तियों,
विधियों और आरेखों से बना पाठ, जहाँ भी उस स्तर पर तर्क पहुँच में हो वहाँ सावधानी
से रखा गया तर्क, और उसके बाद सोच-समझकर चुने गए अभ्यास। जो परिणाम पूरी
व्युत्पत्ति के बिना कहे गए हैं, उन्हें स्पष्ट रूप से स्वीकृत लिखा गया है। सभी
अभ्यासों के पूरे हल पुस्तक के अंत में इकट्ठे हैं। रासायनिक समीकरण हमेशा
संतुलित होते हैं, और किसी वास्तविक पदार्थ के बारे में दी गई हर संख्या किसी
संदर्भ सारणी से ली गई है।

पाठ्यक्रम चार पुस्तकों में है। पहली पुस्तक विद्यालय के सभी बारह वर्षों को एक ही
खंड में समेटती है और किसी भी धारणा को दो बार नहीं पढ़ाती: हर विचार एक ही बार
पढ़ाया जाता है, उस वर्ष में जहाँ वह पहली बार ठीक बैठता है, और बाद का जो अध्याय
उसका उपयोग करता है, वह आगे बढ़ने से पहले एक छोटे-से अनुस्मारक,
\emph{पहले से ज्ञात}, से शुरू होता है। बाकी तीन पुस्तकें विश्वविद्यालय के तीन
वर्ष हैं। पहली पुस्तक में विद्यालय का हर वर्ष एक भाग है, जो अध्यायों में बँटा
है; विश्वविद्यालय की हर पुस्तक एक वर्ष है, जो अध्यायों में बँटी है। वर्ष जितना
पहले का, पाठक उतना ही छोटा: ज़्यादा चित्र, ज़्यादा विस्तार से समझाए गए चरण,
और ज़्यादा सरल अभ्यास।

\section*{इस पुस्तक को कैसे पढ़ें}

कथनों पर हर अध्याय के भीतर क्रमांक दिए गए हैं और वे रंगों से पहचाने जाते हैं:
परिभाषाएँ नीली, प्रमेय लाल, प्रतिज्ञप्तियाँ और उनसे जुड़े कथन नारंगी, और
विधियाँ --- यानी आम कामों के छोटे नुस्ख़े --- हरी। अभ्यास सीधे उपयोग वाले
($\star$) से लेकर कठिन समस्याओं ($\star\star\star$) तक श्रेणियों में बँटे हैं।
फ़्रेम वाले खाने उन बातों के लिए हैं जो तर्क का हिस्सा नहीं हैं: किसी पिछले
अध्याय में सीखी गई बात का अनुस्मारक, प्रयोगशाला में कोई प्रक्रिया कैसे की जाती
है, इतिहास का एक पन्ना, या किसी उत्पाद के ख़तरे।

\section*{सुरक्षा}

रसायन विज्ञान असली पदार्थों से सीखा जाता है, और उनमें से कुछ ख़तरनाक होते हैं।
इस पुस्तक में बताए गए प्रयोग विद्यालय या विश्वविद्यालय की प्रयोगशाला में, किसी
शिक्षक द्वारा या उनके साथ, प्रयोगशाला के सुरक्षा उपकरणों के साथ किए जाते हैं।
इनमें से कोई भी प्रयोग घर पर करने के लिए नहीं है।

\section*{योगदान}

इस पुस्तक का स्रोत एक सार्वजनिक git भंडार में है:
\begin{center}
\url{https://github.com/bvirrion/one-chemistry-book}
\end{center}
योगदान --- नए अध्याय, सुधार, बेहतर व्याख्याएँ, अतिरिक्त अभ्यास ---
सबका स्वागत है; इसके लिए भंडार के मूल में रखी \texttt{CONTRIBUTING.md}
फ़ाइल देखें।
